\documentclass[11pt]{article}
\usepackage{amsmath,amssymb}
\usepackage[margin=1in]{geometry}
\usepackage{hyperref}
\emergencystretch=3em

\title{Programme P: rewriting \S3 of the grammar paper around the Taylor tree}
\author{Claude, with Daniel Murfet}
\date{2026-09-09}

\begin{document}
\maketitle
\sloppy

\begin{abstract}
The user asked for the population theorem of \S3, flagged by its authors as out of date, to be
redesigned around the completed \S4 machinery and formalised. Astra \#37 supplied the design:
Theorem A ``population expansions from admissible normal-form data'' (chart expansion at zero
empirical process, first-candidate exponent with the face-functional coefficient, monomial shift,
tangential integration, finite assembly with the min-exponent/max-log rule) and Corollary B
(leading expectations, three regimes), with wall-crossing demoted to a remark and a 20-unit cap.
Sixteen theorem units (290--305) and five reviews (v32--v36, all PASS) later, the seabed holds the
rewritten leading-order \S3 as conditional Lean theorems, two regression examples that refute the old
statement within one chart, and a staging note with the replacement statements and an errata list.
\end{abstract}

\section{What changed in the mathematics}
The old theorem read the leading coefficient at the deepest corner and called the first candidate
exponent the leading exponent. Both are wrong in general: the coefficient of $N^{-\lambda}(\log
N)^{m-1}$ is the face functional $\Gamma(\lambda)\beta^{-\lambda}/(m-1)!\prod_J(2k_i)^{-1}\int\eta(P_Ju)
\prod_{i\notin J}u_i^{h_i-2k_i\lambda}du$, and it can vanish for a signed amplitude, in which case the
true leading term may sit at the same exponent with a smaller log degree. The formal proof route was
the one review v32 prescribed: isolate the exponent below the cutoff $\lambda+1/Q$ using candidate
support, show the normalised remainder vanishes, transfer Headline VIII's limit to the polynomial in
$\log N$, and read off coefficients by a polynomial growth-uniqueness lemma proved through Mathlib's
polynomial asymptotics. Everything downstream (shift, positivity, tangential integration, assembly,
quotients, the energy observable) reused frozen interfaces: the coefficient-family Taylor tree,
Headline VIII, the Programme S remainder majorants, and the closing corollaries of Programme S.

\section{Lessons}
\begin{itemize}
\item \emph{Zero noise is a bridge, not a definition.} The stochastic-data interface represents a
chart datum in $\ell^1$; the population case is the condition ``phase coordinates vanish''. The link
between the canonical coefficient and the Taylor-tree coefficient at $b=1$ had to be a Lean lemma
(\texttt{boxCoeff\_one}), as review v34 insisted; it was three lines once the prefactors were seen to
be $1$ and $\log1=0$ to kill the binomial re-expansion.
\item \emph{Hypothesis patterns beat finset minima.} The global first candidate $(\mu_*,m_*)$ was
supplied as ``$\mu_*\le\lambda_I$ for all $I$ and attained; $m_*\ge m_I$ on the tie and attained''.
No \texttt{Finset.min'} bookkeeping, and the theorems remain usable when the user knows the answer.
\item \emph{Name clashes surface only in the root build.} \texttt{clampCube} already existed; the
single-module build passed and the full build failed after the push. The exit-gated full build before
push is not optional.
\item \emph{Say ``first candidate''.} Every headline now distinguishes the first candidate from the
leading pair; the two regression examples exist to make the distinction unforgettable.
\end{itemize}

\section{Hand-off}
Report \texttt{population-normal-form-report.pdf}, statements and errata
\texttt{sec3-rewrite-note.md} (both in \texttt{projects/grammar/staging/}); mirror annotations in
\S3 of \texttt{grammar\_lean.tex} under the rule ``label the proved restricted statement and spell out
the restriction''. The geometric content of \S3 (resolution, adapted partition of unity, analytic
admissibility of localised amplitudes) remains external and is the natural next author-side task.
\end{document}
